\begin{theorem}[S-40: Pure tensor rank restrictions without a Brauer-splitting hypothesis]
\label{thm:semireg-pure-tensor-unsplit-rank}
Use the balanced real-quadratic fourfold notation of S-38. Thus
$\Delta_0>1$ is squarefree, $D=\int_X\theta^4/24$, and an actual
perfect object $\mathcal E$ of positive virtual rank $N$ satisfies
\[
 e^{-c_1(\mathcal E)/N}\operatorname{ch}(\mathcal E)=Nv_q.
\]
Suppose $\mathcal E$ has a coherent lift for a finite pure tensor
subgroup $T\subset\widehat X$. Let $p:B\to X$, $r$, $d$, $m=N/(rd)$,
and $\beta$ be its center cover, center degree, Azumaya degree,
reduced rank, and Brauer class as in S-38. Assume
\[
 k=\frac{4DN^2q^2}{|T|}\in\{2,4\}.
 \tag{S40.1}
\]
In particular, this numerical hypothesis holds for a simple, gluable,
fully semiregular lift. Then, without assuming $\beta=0$,
\[
 \boxed{\Delta_0\mid m,\qquad k=2\Longrightarrow2\mid m.}
 \tag{S40.2}
\]

Suppose in addition that $\theta$ is principal, the full maximal
order $\mathcal O_F$ acts integrally, and
$\operatorname{NS}(X)_{\mathbf Q}=F\theta$, as in S-36. If
$\Delta_F$ is the discriminant of $F=\mathbf Q(\sqrt{\Delta_0})$,
put
\[
 d_F=\begin{cases}
 \Delta_F,&\Delta_F\text{ odd},\\
 \Delta_F/2,&\Delta_F\text{ even}.
 \end{cases}
\]
The original virtual rank then satisfies
\[
 \boxed{\begin{aligned}
 k=4&\Longrightarrow d_F\Delta_0\mid N,\\
 k=2&\Longrightarrow 2d_F\operatorname{lcm}(\Delta_0,2)\mid N.
 \end{aligned}}
 \tag{S40.3}
\]
These are necessary conditions; no existence at the indicated
minimal rank is asserted.
\end{theorem}
\begin{proof}
S-38 gives $m\beta=0$. On the Brauer group of a complex abelian
variety the pullback $[a]^*$ is multiplication by $a^2$.
The precise primary statement is Orr--Skorobogatov--Valloni--Zarhin,
\emph{Invariant Brauer group of an abelian variety},
arXiv:2007.05473v3, Section~3, the paragraph immediately following
Remark~3.3, printed p.~6. Its Kummer-sequence input is Stacks
Lemma~59.28.1, tag~03PK. Hence the multiplication isogeny
\[
 f=[m]:B'\longrightarrow B
 \tag{S40.4}
\]
has degree $m^8$ and kills $\beta$. The notation $B'$ denotes the
domain equipped with its pulled-back structures, even though its
underlying abelian variety is again $B$.

Choose a rank-$d$ splitting bundle $S'$ for
$f^*\mathscr A_B$. It has the full exponential character
\[
 \operatorname{ch}(S')=d\exp\!\left(\frac{c_1(S')}{d}\right).
 \tag{S40.5}
\]
Indeed S-38 constructs a further isogeny that makes $\mathscr A_B$
a constant matrix algebra. A connected component of its base
change along $f$ is an isogeny to $B'$ with the same property for
$f^*\mathscr A_B$. On that cover $S'$ is a line bundle tensored
with a constant $d$-dimensional vector space, by Stacks
Lemma~59.62.1. Isogeny pullback is injective on rational
cohomology, giving (S40.5) in every degree.

For each $i$, let $M_i$ be the coherent $\mathscr A_B$-module
corresponding to the $T$-equivariant sheaf
$\mathcal H^i(\mathcal E)$. Ordinary Morita equivalence for the
split algebra expresses
\[
 f^*M_i=S'\otimes F'_i
\]
with $F'_i$ a coherent sheaf on $B'$. Set
$F'=\bigoplus_i F'_i[-i]$. It is an actual perfect object because
$B'$ is smooth, and its virtual rank is $m$. Pulling back the
finite-etale Riemann--Roch identity for $p_*M_i$ gives
\[
 f^*p^*\operatorname{ch}(\mathcal E)
   =r\operatorname{ch}(S')\operatorname{ch}(F').
 \tag{S40.6}
\]
Here the coefficient is the original center degree $r$:
it comes from $p^*p_*=r\operatorname{id}$ on rational cohomology,
before applying $f^*$. It is not the degree of $p\circ f$.
The degree-two component of (S40.6), together with $N=rdm$,
therefore yields
\[
 e^{-c_1(F')/m}\operatorname{ch}(F')
    =m f^*p^*v_q.
 \tag{S40.7}
\]

The integral polarization $\theta'=f^*p^*\theta$ has volume
$D'=m^8rD=m^8D_B$. The rational real multiplication, its Rosati
symmetry, and its type $(2,2)$ transport through these isogenies.
Apply the integral Chern-class conclusion (S30.3) to $F'$.
Using $|T|=rd^2$, $N=rdm$, and hence
$k=4D_Bm^2q^2$, it gives
\[
 m^4D'q^2
   =m^{12}D_Bq^2
   =\frac{m^{10}k}{4}
   \in\Delta_0\mathbf Z_{>0}.
 \tag{S40.8}
\]
An odd prime dividing $\Delta_0$ but not $m$ would make the
middle expression a unit at that prime, which is impossible.
If $m$ is odd, $k=2$ makes it a half-integer; for $k=4$ it is
odd and is incompatible with $2\mid\Delta_0$. These observations
prove (S40.2). No simplicity, gluability, or semiregularity of
$F'$ is used.

For the principal maximal-order case, S-36 gives
$u=Nq\in d_F\mathbf Z_{>0}$. If $k=4$, (S40.1) gives
$rd^2=|T|=u^2$. As $r$ is an integer, it follows that
$r=s^2$ and $sd=u$ for a positive integer $s$. Thus
\[
 N=rdm=s u m,
\]
and the first divisibility in (S40.3) follows from
$d_F\mid u$ and $\Delta_0\mid m$.
If $k=2$, the same identity gives $rd^2=2u^2$.
Consequently $r=2s^2$ and $sd=u$ for a positive integer $s$,
so
\[
 N=2s u m.
\]
Now $\operatorname{lcm}(\Delta_0,2)\mid m$ by (S40.2), proving
the second assertion of (S40.3).
\end{proof}

\begin{remark}[S-40: Concrete principal-field restrictions]
For $F=\mathbf Q(\sqrt2)$, the two rank divisibilities are
$8\mid N$ when $k=4$ and $16\mid N$ when $k=2$.
For $F=\mathbf Q(\sqrt3)$ they are $18\mid N$ and $72\mid N$;
for $F=\mathbf Q(\sqrt5)$ they are $25\mid N$ and $100\mid N$.
All retain the principal, maximal-order, and minimal divisor-space
hypotheses in (S40.3), and concern a pure tensor action.
\end{remark}
